\documentclass[a4paper, 11pt]{report}

\usepackage{a4wide,cite}
\usepackage[utf8]{inputenc}
\usepackage{amsmath, amsthm, amssymb, units, dsfont, algorithmicx, algpseudocode}   % dsfont pro indikátor jevu
\usepackage[pdftex,
						pdfauthor={Jakub Klemsa},
						pdftitle={Models of self-assembling DNA nanostructures, Jakub Klemsa},
						pdfsubject={},
						pdfkeywords={Wang tiles, aTAM, Turing universality},
						pdfproducer={Jakub Klemsa},
						pdfcreator=dvipdf]{hyperref}

\usepackage{graphicx,color,tabularx,float}
\usepackage{fancyhdr}
\pagestyle{fancy}
\fancyhf{}
\fancyhead[L]{\slshape \rightmark} % \leftmark .. kapitola
\fancyhead[R]{\thepage}
\oddsidemargin=10mm   % jednostranný tisk !!!

\let\tmpitem\item
\renewcommand{\item}{\setlength{\itemsep}{2pt}\tmpitem}   % míň roztahovačný seznamy
\renewcommand{\labelitemi}{$\circ$}
%~ \def\refname{Li­teratura}   % záhadně nefunguje ???
%~ \frenchspacing   % údajně česká typografická pravidla

\theoremstyle{definition}
\newtheorem{example}{Example}[chapter]
%~ \newtheorem{definition}{Definice}[section]
%~ \newtheorem{predpoklad}{Předpoklad}[section]
%~ \newtheorem{tvrzeni}{Tvrzení}[section]
%~ \newtheorem{pozorovani}[tvrzeni]{Pozorování}
%~ \newtheorem{lemma}[tvrzeni]{Lemma}
%~ \newtheorem{dusledek}[tvrzeni]{Důsledek}
%~ \newtheorem{algoritmus}{Algoritmus}[chapter]

%~ \theoremstyle{remark} 
%~ \newtheorem{pozn}{Poznámka}[chapter]

\DeclareMathOperator*{\E}{E\,}   % operátor* bude umět používat \limit :)
\DeclareMathOperator{\de}{d\!}   % v integrálech
\renewcommand{\P}{\mathbb{P\,}}
\newcommand{\R}{\mathbb{R}}
\newcommand{\Z}{\mathbb{Z}}
\newcommand{\N}{\mathbb{N}}
\newcommand{\indicator}{\mathds{1}}
\newcommand{\powerset}[1]{\mathcal{P} ( #1 )}
\newcommand{\adob}[2]{#1 , \ldots , #2}

%%% VYHLEDAT A OPRAVIT !!! %%%%%%%%%%%%%%%%%%%%%%%%%%%%%%%%%%%%%%%%%%%%%
\let\haluz\relax
\haluz % program vlna nakonec!
% vlna !!
% taky opravit N-1 na N\!-\!1
% zkontrolovat \min \; ...
%%%%%%%%%%%%%%%%%%%%%%%%%%%%%%%%%%%%%%%%%%%%%%%%%%%%%%%%%%%%%%%%%%%%%%%%

%%% k vyplnění ZAČÁTEK -------------------------------------------------
\newcommand{\cvut}{ČESKÉ VYSOKÉ UČENÍ TECHNICKÉ V~PRAZE}
\newcommand{\fjfi}{Fakulta jaderná a fyzikálně inženýrská}
\newcommand{\km}{Katedra matematiky}
\newcommand{\obor}{Inženýrská informatika}
\newcommand{\minf}{Matematická informatika}
\newcommand{\TYPPRACE}{VÝZKUMNÝ ÚKOL}
\newcommand{\TypPrace}{Výzkumný úkol}
\newcommand{\mePrace}{mého výzkumného úkolu}
\newcommand{\rodPrace}{ho}

\newcommand{\nazevcz}{Modely sebeskládajících DNA nanostruktur}
\newcommand{\nazeven}{Models of self-assembling DNA nanostructures}
\newcommand{\autor}{Jakub Klemsa}
\newcommand{\vedouci}{Ing. Štěpán Starosta, Ph.D.}
\newcommand{\vedoucimu}{Ing. Štěpánu Starostovi, Ph.D.}
\newcommand{\pracovisteVed}{}
\newcommand{\konzultant}{---}
\newcommand{\rok}{2014}
\newcommand{\akadRok}{2013/2014}
\newcommand{\misto}{Praha}
\newcommand{\vMiste}{Praze}

\newcommand{\klicova}{}
\newcommand{\keyword}{}
\newcommand{\abstrCZ}{Bla bla abstrakt CZ.}
\newcommand{\abstrEN}{Bla bla abstrakt EN.}
%%% k vyplnění KONEC ---------------------------------------------------


\begin{document}


\pagenumbering{roman}
\begin{titlepage}

%%% 1. úvodní strana
\thispagestyle{empty}
\begin{center}
	{\Large \cvut \\[12pt] \fjfi \\[260pt]}
	{\Huge \TYPPRACE}
\end{center}
\vfill
{
	\Large \misto, \rok \hfill \autor
}
\newpage


%%% 2. úvodní strana
\thispagestyle{empty}
\begin{center}
	{\Large \cvut \\[10pt] \fjfi \\[10pt] \km \\[40pt]}
	\includegraphics[height=75pt]{lev.eps} \\[100pt]
	
	{\LARGE \TYPPRACE \\[60pt]}
	
	{\Large\bf \nazevcz \\[30pt]}
	
	{\Large\bf \nazeven}
\end{center}
\vfill
{
	\Large
	\begin{tabular}{ll}
	Vypracoval: & \autor\\[3pt]
	Školitel: & \vedouci\\[3pt]
	Akademický rok: & \akadRok
	\end{tabular}
}
\newpage


%%% sem přijde ofic. zadání
\thispagestyle{empty}
\Large
Na toto místo přijde svázat \textbf{zadání \mePrace}!

\vspace{4mm}
V~jednom z~výtisků musí být \textbf{originál} zadání, v~ostatních kopie.
\normalsize
\newpage


%%% prohlášení
\thispagestyle{empty}
~
\vfill
\noindent\textbf{Čestné prohlášení}
\vspace{0.5cm}

\noindent
Prohlašuji na tomto místě, že jsem předloženou práci vypracoval samostatně a že jsem uvedl veškerou použitou literaturu.
\vspace{1.5cm}

\noindent
\vspace{5mm}V \vMiste dne \today\hfill
	\begin{tabular}{c}
	. . . . . . . . . . . . . . . . . . . . . . .\\
	\autor
	\end{tabular}
\newpage


%%% poděkování
\thispagestyle{empty}
~
\vfill
\noindent\textbf{Poděkování}
\vspace{0.5cm}

\noindent
Děkuji \vedoucimu, za vedení \mePrace a za podnětné návrhy, které \rodPrace ~obohatily.

\begin{flushright}
\autor
\end{flushright}
\newpage


%%% abstrakt atp.
\thispagestyle{empty}

\begin{tabularx}{\textwidth}{lX}
  {\em Název práce:} & \bf \nazevcz \\[4mm]
  {\em Autor:} & \autor \\[4mm]
  {\em Obor:} & \obor \\[4mm]
  {\em Zaměření:} & \minf \\[4mm]
  {\em Druh práce:} & \TypPrace \\[4mm]
  {\em Vedoucí práce:} & \vedouci, \pracovisteVed \\[4mm]
  {\em Konzultant:} & \konzultant \\[4mm]
  {\em Abstrakt:} & \abstrCZ \\[4mm]
  {\em Klíčová slova:} & \klicova \\[20mm]

  {\em Title:} & \bf \nazeven \\[4mm]
  {\em Author:} & \autor \\[4mm]
  {\em Abstract:} & \abstrEN \\[4mm]
  {\em Key words:} & \keyword
\end{tabularx}
\newpage


%%% obsah
\tableofcontents
\thispagestyle{empty}

\end{titlepage}
\pagenumbering{arabic}


%%%%%%%%%%%%%%%%%%%%%%%%%%%%%%%%%%%%%%%%%%%%%%%%%%%%%%%%%%%%%%%%%%%%%%%%
%     Z A Č Á T E K   P R Á C E                                        %
%%%%%%%%%%%%%%%%%%%%%%%%%%%%%%%%%%%%%%%%%%%%%%%%%%%%%%%%%%%%%%%%%%%%%%%%


%%% ÚVODEM
\cleardoublepage\phantomsection   % cleardoublepage a phantomsection udělají že odkaz z obsahu vede NAD nadpis, jinak vede POD
\addcontentsline{toc}{chapter}{Prolog}
\chapter*{Prolog}

	1959: Feynman's visionary talk, \cite{feynman}; 
	
	The ground-breaking work was carried out by Adleman, \cite{adleman94}, who showed that DNA computation is practically feasible. In his experiment, Adleman used special DNA sequences for solving Hamiltonian Path Problem, one of the most typical NP-complete problems.
	
	... Extreme parallelism! But also possibility of errors.

\section*{Work overview}
	
	%%%%%%%%%%%%%%%%%%%%%%%%%%%%%%%%%%%%%%%%%%%%%%%%%%%%%%%%%%%%%%%%%%%%%%
	
	Chapter 1: Intro.\\
	Chapter 2: First of all I will describe models which exploit specific DNA structure.\\
	Chapter 3: Abstract Tile Assembly Model, temperature, 2D vs. 3D.\\
	Positive integers $\N$. \nomenclature{$\N$}{Positive integers.}
	
	% zkusíme Ramseyovo číslo R(5,5) ? :)
	% -> problém "Je R(n,n) větší než něco?" znamená pro odpověď ANO najít graf t.ž. neobsahuje K_n ani anti-K_n .. což je samo o sobě co-NP (neboli to je \exists \forall stroj .. \Sigma_2 jazyk, horní hranice pak je \Pi_2 jazyk)
	% 
	% NP vs co-NP .. co-NP přece nedává ne?

%%% KAPITOLA 1 -- Úvod
\chapter{Overview}
	
	Areas of interest.

%%% KAPITOLA 2 -- Testy
\chapter{Computational power}
	
	How powerful aTAM model is? Proved by Winfree, see \cite{winfree_phd}. Many other results in \cite{cook_temp1}, \cite{stage_assembly}, \cite{square_lb}, \cite{square_ub} \ldots
	
	\begin{figure}[H]
	\begin{center}
		\includegraphics{./plots/tile-1.eps}
		\caption{Tile}
	\end{center}
	\end{figure}

%%% ZÁVĚREM
\cleardoublepage\phantomsection   % cleardoublepage a phantomsection udělají že odkaz z obsahu vede NAD nadpis, jinak vede POD
\addcontentsline{toc}{chapter}{Epilog}
\chapter*{Epilog}
	
	The very last section to be done.

%%% LITERATURA
% \renewcommand\bibname{Literatura}
\cleardoublepage\phantomsection   % cleardoublepage a phantomsection udělají že odkaz z obsahu vede NAD nadpis, jinak vede POD
\addcontentsline{toc}{chapter}{References}

\bibliography{biblio}{}
\bibliographystyle{plain}


%%%%%%%%%%%%%%%%%%%%%%%%%%%%%%%%%%%%%%%%%%%%%%%%%%%%%%%%%%%%%%%%%%%%%%%%
%     K O N E C   P R Á C E                                            %
%%%%%%%%%%%%%%%%%%%%%%%%%%%%%%%%%%%%%%%%%%%%%%%%%%%%%%%%%%%%%%%%%%%%%%%%


\end{document}